% from ega4-2-fr.tex
\oldpage[IV-2]{12}
fidèlement plat, et par suite le foncteur
$\mathscr G_y\mapsto\mathscr G_y\otimes_{\mathcal O_y}\mathscr F_x$ est
fidèle~; comme par ailleurs pour \emph{tout} $x\in f^{-1}(y)$ le foncteur
$\mathscr G_y\mapsto\mathscr G_y\otimes_{\mathcal O_y}\mathscr F_x$ est
exact, on en déduit aussitôt la conclusion.
\end{remark}

\begin{proposition}[2.2.10]
\label{IV.2.2.10-fr}
Soient $f:X\to Y$, $g:Y'\to Y$, $h:Y'\to Z$ trois morphismes de préschémas,
$X'=X\times_Y Y'$, $\mathscr F$ un $\mathcal O_X$-Module quasi-cohérent,
fidèlement plat relativement à $Y$, $\mathscr G'$ un
$\mathcal O_{Y'}$-Module quasi-cohérent. Pour que $\mathscr G'$ soit plat
(resp. fidèlement plat) relativement à $Z$, il faut et il suffit que
$\mathscr F\otimes_Y\mathscr G'$ soit un $\mathcal O_{X'}$-Module plat
(resp. fidèlement plat) relativement à $Z$.
\end{proposition}

On sait déjà que si $\mathscr G'$ est $Z$-plat, il en est de même de
$\mathscr F\otimes_Y\mathscr G'$ (2.1.5). Considérons un changement de base
quelconque $Z''\to Z$ et posons $X''=X'\times_Z Z''$~; si $\mathscr G'$ est
fidèlement plat relativement à $Z$, le foncteur
\begin{equation}
\tag{2.2.10.1}
\label{IV.2.2.10.1-fr}
\mathscr H''\mapsto\mathscr H''\otimes_Z\mathscr G'
\mapsto(\mathscr H''\otimes_Z\mathscr G')\otimes_Y\mathscr F
=\mathscr H''\otimes_Z(\mathscr G'\otimes_Y\mathscr F)
\end{equation}
de la catégorie des $\mathcal O_{Z''}$-Modules quasi-cohérents dans celle des
$\mathcal O_{X''}$-Modules est composé de deux foncteurs exacts et fidèles,
donc est exact et fidèle. Inversement, si ce foncteur composé est exact (resp.
exact et fidèle), il en est de même de
$\mathscr H''\mapsto\mathscr H''\otimes_Z\mathscr G'$, puisque le foncteur
$\mathscr M'\mapsto\mathscr M'\otimes_Y\mathscr F$ (de la catégorie des
$\mathcal O_{Y'}$-Modules quasi-cohérents dans celle des
$\mathcal O_{X'}$-Modules quasi-cohérents) est exact et fidèle par hypothèse.

\begin{corollary}[2.2.11]
\label{IV.2.2.11-fr}
\begin{enumerate}
\item[(i)] Soient $f:X\to Y$, $g:Y'\to Y$ deux morphismes,
$X'=X\times_Y Y'$, $\mathscr F$ un $\mathcal O_X$-Module quasi-cohérent. Si
$\mathscr F$ est fidèlement plat relativement à $Y$,
$\mathscr F'=\mathscr F\otimes_Y\mathcal O_{Y'}$ est fidèlement plat
relativement à $Y'$.
\item[(ii)] Soient $f:X\to Y$, $g:Y\to Z$ deux morphismes, $\mathscr F$ un
$\mathcal O_X$-Module quasi-cohérent, fidèlement plat relativement à $Y$.
Pour que $g$ soit un morphisme fidèlement plat, il faut et il suffit que
$\mathscr F$ soit fidèlement plat relativement à $Z$.
\item[(iii)] Soient $f:X\to Y$, $g:Y\to Z$ deux morphismes, $\mathscr G$ un
$\mathcal O_Y$-Module quasi-cohérent. On suppose le morphisme $f$ fidèlement
plat. Pour que $\mathscr G$ soit plat (resp. fidèlement plat) relativement à
$Z$, il faut et il suffit que $f^*(\mathscr G)$ soit plat (resp. fidèlement
plat) relativement à $Z$.
\item[(iv)] Soient $f:X\to Y$, $g:Y\to Z$ deux morphismes, $x$ un point de
$X$. On suppose $f$ plat au point $x$. Pour que $g\circ f$ soit plat au point
$x$, il faut et il suffit que $g$ soit plat au point $f(x)$.
\end{enumerate}
\end{corollary}

Pour démontrer (i), on applique (2.2.10) en remplaçant $Z$ par $Y'$ et
$\mathscr G'$ par $\mathcal O_{Y'}$. Pour démontrer (ii), on applique
(2.2.10) en prenant pour morphisme $Y'\to Y$ l'identité et remplaçant
$\mathscr G'$ par $\mathcal O_Y$. Pour démontrer (iii), on applique (2.2.10)
en prenant encore l'identité pour $Y'\to Y$ et remplaçant $\mathscr F$ par
$\mathcal O_X$ et $\mathscr G'$ par $\mathscr G$. Enfin (iv) se déduit de
(ii) en remplaçant $X$ par $\operatorname{Spec}(\mathcal O_x)$,
$\mathscr F$ par $\mathcal O_X$, $Y$ par
$\operatorname{Spec}(\mathcal O_{f(x)})$ et $Z$ par
$\operatorname{Spec}(\mathcal O_{g(f(x))})$, et tenant compte de
$(0_{\mathrm I},6.6.2)$.

\begin{corollary}[2.2.12]
\label{IV.2.2.12-fr}
Soient $Y$ un schéma affine, $f:X\to Y$ un morphisme quasi-compact,
$\mathscr F$ un $\mathcal O_X$-Module quasi-cohérent. Si $\mathscr F$ est
fidèlement plat relativement à $f$, il existe un schéma affine $X'$ et un
isomorphisme local surjectif $g:X'\to X$ tels que $g^*(\mathscr F)$ soit
fidèlement plat relativement à $f\circ g$.
\end{corollary}

En effet, $X$ est quasi-compact, donc réunion finie d'ouverts affines $X_i$~;
il suffit de prendre pour $X'$ le schéma affine \emph{somme} des préschémas
induits sur les ouverts $X_i$ et
